
\chapter{Conclusiones}
Un conjunto de datos de calidad es el requisito fundamental para establecer la base de un modelo de aprendizaje profundo (deep learning), 
ya que el rendimiento de dicho modelo depende en gran medida de la calidad, cantidad y relevancia del conjunto de datos. Por lo tanto,
el primer paso crucial de cualquier aprendizaje profundo comienza con la adquisición de imágenes.

El proceso de adquisición de datos no solo implica capturar la información, sino asegurar que esta sea representativa del entorno real donde 
se desplegará el modelo, lo cual es el primer filtro para evitar sesgos y errores en etapas posteriores de entrenamiento \cite{Khan2023}.